% arXiv-style template (single-column, similar to ACL/arXiv preprints)
\documentclass[11pt]{article}

\usepackage[utf8]{inputenc}
\usepackage[T1]{fontenc}
\usepackage{microtype}
\usepackage[margin=1in]{geometry}
\usepackage{graphicx}
\usepackage{booktabs}
\usepackage{amsmath,amssymb}
\usepackage{hyperref}
\usepackage{url}
\usepackage{natbib}
\usepackage{xcolor}
\usepackage{authblk}

\hypersetup{
  colorlinks=true,
  linkcolor=blue!60!black,
  citecolor=blue!60!black,
  urlcolor=blue!60!black,
}

\title{Benchmarking Classical, Embedding-Based, and Transformer Approaches\\
       for Hate Speech Detection in Online Comments}

\author[1]{Shahnawaz Khan}
\affil[1]{QA Pulse by SK \quad \texttt{shahnawaz.jrw@gmail.com}}

\date{\today}

\begin{document}
\maketitle

\begin{abstract}
We address the problem of hate speech detection in online user comments. Hate
speech --- abusive speech targeting specific group characteristics such as
ethnicity, religion, or gender --- is an important problem plaguing websites
that allow users to leave feedback, with a negative impact on online business
and overall user experience. We benchmark three approaches that span two
decades of progress in text classification: (i) a TF-IDF + logistic regression
baseline, (ii) Paragraph2Vec (Doc2Vec) comment embeddings followed by a linear
classifier, replicating \citet{djuric2015hate}, and (iii) fine-tuned
DistilBERT. We evaluate on two public datasets, \citet{davidson2017automated}
and HateXplain \citep{mathew2021hatexplain}, and report a cross-dataset
generalization study. Our results confirm that distributed representations
address the high-dimensionality and sparsity issues of bag-of-words
approaches, while contextual transformer models yield further consistent gains
in F$_1$ on the hate class. Code and configurations are released at
\url{https://github.com/ShahnawazKakarh/hate-speech-detection}.
\end{abstract}

\section{Introduction}

In the age of ever-increasing volume and complexity of the internet, millions
of users have unrestricted access to vast amounts of content that allows for
privileges unimaginable several decades ago, such as access to knowledge
bases or latest news within just a few clicks. However, due to the
non-restrictive nature of the medium and, in certain jurisdictions, legal
protection of free speech that also includes hate speech, some users misuse
the medium to promote offensive and hateful language. This mars the
experience of regular users, affects the business of online companies, and
may even have severe real-life consequences.

To mitigate these detrimental effects, many companies (including Yahoo,
Facebook, and YouTube) prohibit hate speech on websites they own and operate,
and implement algorithmic solutions to discern hateful content. However, the
scale and multifacetedness of the task render it a difficult endeavor, and
hate speech remains a problem in online user comments. Despite the prevalence
and large impact of online hate speech, there exist comparatively few
unified, reproducible benchmarks that compare classical, embedding-based, and
modern transformer approaches on shared data with shared evaluation
protocols. This paper contributes such a benchmark.

\paragraph{Contributions.} We make the following contributions:
\begin{enumerate}
  \item A unified, configuration-driven training pipeline covering three
        canonical approaches: TF-IDF + LR, Doc2Vec + LR, and DistilBERT
        fine-tuning.
  \item Evaluation on two public hate-speech datasets
        (\citealp{davidson2017automated}; HateXplain) with consistent
        preprocessing and metrics.
  \item A cross-dataset generalization study quantifying how well each model
        transfers from one corpus to another.
  \item Open-source release of code, configurations, and prepared dataset
        splits to support reproducibility.
\end{enumerate}

\section{Background and Related Work}

\paragraph{Real-world impact of online hate speech.}
Recent cases highlight the impact of hurtful language in online communities
as well as on major corporations. For example, in 2013 Facebook came under
fire for hosting pages that were hateful against women; a petition amassed
over 200{,}000 supporters within days, and several major advertisers
threatened to pull their ads. At the individual level, when actor Robin
Williams passed away, his daughter Zelda was bullied on Twitter and
Instagram and eventually deleted all of her online accounts, prompting
Twitter to revise its hate-speech guidelines. Any platform hosting
user-generated content must therefore contend with a moderation problem.

\paragraph{Detection approaches.}
\citet{djuric2015hate} proposed a two-step method: first learn distributed
low-dimensional representations of comments using \texttt{paragraph2vec} with
the continuous-bag-of-words (CBOW) variant
\citep{le2014distributed,mikolov2013distributed}, then train a binary
classifier on top of these embeddings. The approach addresses the
high-dimensionality and sparsity issues that arise with bag-of-words and
$n$-gram representations.
\citet{nobata2016abusive} developed a supervised classification pipeline
combining $n$-grams, linguistic, syntactic, and distributional semantic
features, evaluated longitudinally on Yahoo Finance and News comments.
\citet{chen2012detecting} introduced the Lexical Syntactic Feature (LSF)
framework, incorporating user-level features for offensive content detection
in social media.
\citet{xiang2012detecting} explored bootstrapping for vulgar language
detection on Twitter. More recently,
\citet{davidson2017automated} released a Twitter dataset distinguishing
hate speech from merely offensive language, and \citet{mathew2021hatexplain}
released HateXplain, which augments labels with rationale annotations.
Transformer-based models such as BERT \citep{devlin2019bert} and its
distilled variant DistilBERT \citep{sanh2019distilbert} have since become
strong baselines for text classification, including hate-speech detection.

\paragraph{Terminology.} Throughout we adopt the following working
definitions, consistent with \citet{nobata2016abusive}:
\textit{hate speech} attacks or demeans a group based on protected
attributes (race, religion, gender, etc.); \textit{derogatory} language
attacks individuals or groups without rising to hate speech;
\textit{profanity} contains sexual remarks or vulgarity.

\section{Datasets}

We use two public, English-language hate-speech corpora.

\paragraph{Davidson 2017.} The corpus of \citet{davidson2017automated}
contains roughly $24{,}783$ tweets labeled into three categories:
\textit{hate}, \textit{offensive}, and \textit{neither}. We use the official
class column directly and binarize as \textsc{hate} vs.\ \textsc{non-hate}.

\paragraph{HateXplain.} \citet{mathew2021hatexplain} released roughly
$20{,}148$ posts from Twitter and Gab, each annotated by three workers into
\textit{hatespeech}, \textit{offensive}, or \textit{normal}. We use majority
vote and the official train/val/test split, then binarize.

For both datasets we apply identical preprocessing: URL/mention removal,
hashtag-word retention, emoji-to-token conversion, lowercasing, and
whitespace normalization.

\section{Methods}

\subsection{TF-IDF + Logistic Regression}
A canonical baseline: unigram + bigram TF-IDF features
($\max=50{,}000$, $\min\_df=2$, sublinear TF) fed to an $L_2$-regularized
logistic regression with class-balanced weights.

\subsection{Doc2Vec + Logistic Regression}
Following \citet{djuric2015hate}, we jointly learn comment and word
embeddings using gensim's PV-DBOW variant of \texttt{Doc2Vec}
\citep{le2014distributed} with $d=200$, window $5$, and $20$ epochs; we
additionally train word vectors via \texttt{dbow\_words=1}. A logistic
regression is then trained on the resulting comment vectors. At inference
time, vectors for unseen comments are obtained via the standard
\emph{folding-in} (\texttt{infer\_vector}) procedure.

\subsection{DistilBERT}
We fine-tune \texttt{distilbert-base-uncased} \citep{sanh2019distilbert} for
binary sequence classification with maximum sequence length $128$, learning
rate $2{\times}10^{-5}$, batch size $32$, $3$ epochs, weight decay $0.01$,
and a $10\%$ linear warm-up. The best checkpoint by validation macro-F$_1$
is used for test-set evaluation.

\section{Experiments}

\subsection{Setup}
All experiments use seed $42$ and the same 80/10/10 stratified train/val/test
split for Davidson; HateXplain uses the official split. We report accuracy,
precision and recall on the hate class, binary F$_1$ on the hate class,
macro-F$_1$, and ROC-AUC.

\subsection{Main results}

% placeholder; overwritten by scripts/make_results_table.py once metrics exist
\begin{tabular}{llcccccc}
\toprule
Model & Dataset & Acc & P & R & F1 (hate) & F1 (macro) & AUC \\
\midrule
TF-IDF + LR  & davidson    & -- & -- & -- & -- & -- & -- \\
TF-IDF + LR  & hatexplain  & -- & -- & -- & -- & -- & -- \\
Doc2Vec + LR & davidson    & -- & -- & -- & -- & -- & -- \\
Doc2Vec + LR & hatexplain  & -- & -- & -- & -- & -- & -- \\
DistilBERT   & davidson    & -- & -- & -- & -- & -- & -- \\
DistilBERT   & hatexplain  & -- & -- & -- & -- & -- & -- \\
\bottomrule
\end{tabular}


\subsection{Cross-dataset generalization}
We additionally evaluate each Davidson-trained model on the HateXplain test
split, and vice versa. Results are written to
\texttt{artifacts/cross\_dataset.json} after running
\texttt{scripts/cross\_dataset\_eval.py}.

\section{Discussion}

\paragraph{What the embeddings buy you.}
Replacing sparse $n$-gram features with $200$-dimensional Doc2Vec embeddings
addresses the dimensionality and sparsity issues highlighted by
\citet{djuric2015hate}, and we expect the embedding space to place
swearwords, slurs, and their obfuscated variants in the same neighborhood,
consistent with the qualitative findings in that paper.

\paragraph{What transformers add.}
DistilBERT brings contextual modeling and subword tokenization, which we
expect to help in cases where context disambiguates surface offensiveness
(e.g., reclaimed slurs) and where adversarial obfuscation defeats
whitespace-tokenized models.

\paragraph{Limitations.} We binarize a three-way label; the hate-versus-
offensive boundary is itself contested. We do not address class imbalance
beyond class weighting, and our evaluation is English-only.

\section{Conclusion and Future Work}

As the volume of online user-generated content grows, accurate automated
methods to flag abusive language are of paramount importance. We provide a
unified, reproducible benchmark of three canonical approaches across two
public datasets. Future work includes multilingual extension, target-group
classification using HateXplain's rationale annotations, and longitudinal
evaluation against language drift.

\bibliographystyle{plainnat}
\bibliography{references}

\end{document}
